\documentclass[11pt]{amsart}
\usepackage[margin=1in]{geometry}
\usepackage{amsmath,amssymb,amsthm}
\usepackage{graphicx}
\usepackage[hidelinks]{hyperref}

\newtheorem{theorem}{Theorem}
\newtheorem{lemma}{Lemma}
\newtheorem{definition}{Definition}
\newtheorem{remark}{Remark}

\newcommand{\A}{\mathcal A}
\newcommand{\K}{\mathbb K}
\newcommand{\Lin}{\mathcal L}
\newcommand{\NA}{\operatorname{NA}}
\newcommand{\Id}{\operatorname{Id}}
\newcommand{\SOT}{\operatorname{SOT}}
\newcommand{\WOT}{\operatorname{WOT}}
\newcommand{\Ball}{B}

\title{A Separably Determined Ideal Obstruction for Universal Norm Attainment}
\author{Agent lane 07}
\date{June 14, 2026}

\begin{document}
\maketitle

\begin{abstract}
We record a partial extension of Dantas--Jung--Martinez-Cervantes'
criterion for the existence of non-norm-attaining operators.  If \(\A\) is a
norm-closed operator ideal separably determined in the domain variable, \(E\)
is reflexive, and the pair \((E,F)\) has the bounded \(\A\)-approximation
property, then universal norm-attainment for operators \(E\to F\) forces
\(\Lin(E,F)=\A(E,F)\).  Thus every proper such ideal yields a
non-norm-attaining operator under bounded \(\A\)-approximation.  The result
recovers the compact-operator case of arXiv:2102.06452 when \(\A=\mathcal K\).
\end{abstract}

\section{Source Problem}

The source paper is:

\begin{quote}
Sheldon Dantas, Mingu Jung, Gonzalo Martinez-Cervantes,
\emph{On the existence of non-norm-attaining operators}, arXiv:2102.06452.
\end{quote}

The paper quotes the following problem of Ostrovskii and the related theorem.

\begin{center}
\includegraphics[width=\textwidth]{figures/open_problem_crop.png}
\end{center}

In words, Ostrovskii asks whether there are infinite-dimensional Banach spaces
on which every operator attains its norm.  Dantas--Jung--Martinez-Cervantes
prove, among other things, that if such a space \(E\) exists, then \(E\) has no
bounded \(\A\)-approximation property for any norm-closed operator ideal
\(\A\) not containing \(\Id_E\) \cite[Theorem 3.8]{DJM}.

The result below is a pair-level version for separably determined ideals.  It
does not solve Ostrovskii's problem; it enlarges the obstruction from compact
operators to a class of closed ideals for pairs \((E,F)\).

\section{Definitions}

\begin{definition}
Let \(\A\) be a norm-closed operator ideal.  We say that \(\A\) is
\emph{separably determined in the domain variable} if for every operator
\(T\in \Lin(X,Y)\),
\[
T\in \A(X,Y)
\quad\Longleftrightarrow\quad
T|_S\in \A(S,Y)\text{ for every separable subspace }S\subset X .
\]
\end{definition}

\begin{definition}
The pair \((E,F)\) has the \emph{bounded \(\A\)-approximation property} if
there is \(\lambda\geq 1\) such that every \(T\in\Lin(E,F)\) belongs to
\[
\lambda\,\overline{\Ball_{\A(E,F)}}^{\,\tau_c},
\]
where \(\tau_c\) is the topology of compact convergence and
\(\Ball_{\A(E,F)}=\{S\in\A(E,F):\|S\|\leq 1\}\).
\end{definition}

\section{Proof Intuition}

Dantas--Jung--Martinez-Cervantes show that if a norm-closed convex subset of
\(\Lin(E,F)\) fails to be sequentially closed in the strong operator topology,
then the pair \((E,F)\) has their James property, and hence there is a
non-norm-attaining operator.  Thus, under universal norm-attainment, closed
convex sets of operators should be strongly sequentially closed.

For the compact ideal, their paper converts SOT non-closure of the compact
unit ball into SOT sequential non-closure by passing to a complemented
separable subspace of the reflexive domain.  The same mechanism only uses the
fact that compactness is detected on separable subspaces.  Therefore it works
for any closed ideal \(\A\) with the same separable-determination property.
Once the ideal unit ball is SOT closed, bounded \(\A\)-approximation leaves no
room for operators outside \(\A(E,F)\).

\section{Main Result}

\begin{lemma}\label{lem:sot}
Let \(E\) be reflexive, let \(F\) be a Banach space, and let \(\A\) be a
norm-closed operator ideal separably determined in the domain variable.  If
\(\Ball_{\A(E,F)}\) is not SOT closed in \(\Lin(E,F)\), then
\(\Ball_{\A(E,F)}\) is not SOT sequentially closed.
\end{lemma}

\begin{proof}
Choose \(T\in \overline{\Ball_{\A(E,F)}}^{\SOT}\setminus
\Ball_{\A(E,F)}\).  Since \(T\) is a pointwise limit of contractions, \(\|T\|\leq
1\).  Hence \(T\notin \A(E,F)\).

Because \(\A\) is separably determined, there is a separable subspace
\(S\subset E\) such that \(T|_S\notin \A(S,F)\).  Reflexive spaces have the
separable complementation property \cite[Lemma 3.5 and its proof]{DJM}; hence
there is a complemented separable subspace \(Z\subset E\) with \(S\subset Z\).
Let \(P:E\to Z\) be a bounded projection.

Since \(T\) lies in the SOT closure of \(\Ball_{\A(E,F)}\), the restriction
\(T|_Z\) lies in the SOT closure of
\[
\{R|_Z:R\in \Ball_{\A(E,F)}\}\subseteq \Ball_{\A(Z,F)}.
\]
The space \(Z\) is separable, so we can choose a sequence
\((R_n)\subset \Ball_{\A(E,F)}\) such that \(R_n z\to Tz\) for every \(z\in Z\).
Let \(C=\max\{1,\|P\|\}\).  Then
\[
S_n:=C^{-1}R_nP\in \Ball_{\A(E,F)}
\]
and \(S_n\to C^{-1}TP\) in the SOT on \(\Lin(E,F)\).

The limit \(C^{-1}TP\) does not belong to \(\A(E,F)\).  Indeed, if it did, then
restricting to \(Z\) would give \(C^{-1}T|_Z\in\A(Z,F)\), and then restricting
again to \(S\) would give \(T|_S\in\A(S,F)\), a contradiction.  Thus
\(\Ball_{\A(E,F)}\) is not SOT sequentially closed.
\end{proof}

\begin{theorem}\label{thm:main}
Let \(E\) be reflexive and \(F\) be a Banach space.  Let \(\A\) be a
norm-closed operator ideal separably determined in the domain variable.  If
every operator from \(E\) into \(F\) attains its norm, then
\(\Ball_{\A(E,F)}\) is SOT closed.

Consequently, if the pair \((E,F)\) has the bounded \(\A\)-approximation
property and every operator \(E\to F\) attains its norm, then
\[
\Lin(E,F)=\A(E,F).
\]
Equivalently, if \(\A(E,F)\neq \Lin(E,F)\) and \((E,F)\) has the bounded
\(\A\)-approximation property, then \(\Lin(E,F)\) contains a non-norm-attaining
operator.
\end{theorem}

\begin{proof}
Assume that every operator \(E\to F\) attains its norm.  If
\(\Ball_{\A(E,F)}\) were not SOT closed, Lemma \ref{lem:sot} would imply that it
is not SOT sequentially closed.  The unit ball \(\Ball_{\A(E,F)}\) is norm
closed and convex.  By \cite[Lemma 3.4]{DJM}, any norm-closed convex subset of
\(\Lin(E,F)\) that is not SOT sequentially closed gives the James property for
the pair \((E,F)\).  By \cite[Theorem A]{DJM}, the James property yields a
non-norm-attaining operator \(E\to F\), a contradiction.  Therefore
\(\Ball_{\A(E,F)}\) is SOT closed.

Now assume that \((E,F)\) has the bounded \(\A\)-approximation property.  Choose
\(\lambda\geq 1\) such that
\[
\Lin(E,F)\subseteq \lambda\,\overline{\Ball_{\A(E,F)}}^{\,\tau_c}.
\]
Since the topology of compact convergence is stronger than the SOT,
\[
\overline{\Ball_{\A(E,F)}}^{\,\tau_c}
\subseteq
\overline{\Ball_{\A(E,F)}}^{\,\SOT}
=\Ball_{\A(E,F)}.
\]
Thus every \(T\in\Lin(E,F)\) belongs to \(\lambda\Ball_{\A(E,F)}\subseteq
\A(E,F)\).  Hence \(\Lin(E,F)=\A(E,F)\).  The final equivalent formulation is
the contrapositive.
\end{proof}

\section{Verification and Scope}

The proof is analytic and uses no computation.  The key verifier checks are:

\begin{itemize}
\item the ideal \(\A\) must be norm closed and separably determined in the
domain variable;
\item the domain \(E\) is reflexive so that the separable complementation
argument used in \cite[Lemma 3.5]{DJM} applies;
\item the approximation hypothesis is bounded and is stated for the pair
\((E,F)\).
\end{itemize}

This does not solve the full Ostrovskii problem.  It says that any positive
example must avoid bounded \(\A\)-approximation for every proper
domain-separably-determined closed ideal \(\A\) in the pair-level sense above.
For \(\A=\mathcal K\), the theorem recovers the compact-operator obstruction
from the source paper.

\section{Novelty Check}

Cheap run indexes had no existing packet for arXiv:2102.06452.  Local searches
for ``separably determined'', ``bounded A-approximation property'', ``operator
ideal'' with ``norm attaining'', and ``James property'' found the source paper
and no existing run result with this theorem.  Web searches for exact phrases
around separably determined operator ideals and non-norm-attaining operators
found no separate later answer.

The arXiv source file of \cite{DJM}, in comments not rendered in the published
PDF, mentions that separably determined ideals should allow a similar
SOT-closure argument.  For that reason this packet should be read as a precise
reviewable extension of the source method, with modest novelty, rather than a
settlement of the original open problem.

\begin{thebibliography}{9}
\bibitem{DJM}
S. Dantas, M. Jung, and G. Martinez-Cervantes,
\emph{On the existence of non-norm-attaining operators},
arXiv:2102.06452.
\end{thebibliography}

\end{document}
